\documentclass[12pt]{amsart}
\usepackage{amsmath,amssymb, bm, amsthm,graphicx}
\usepackage{xcolor}
\usepackage[T2A]{fontenc}
\usepackage[utf8]{inputenc} % исходные файлы в WIN кодировке
\usepackage[russian]{babel}
\usepackage{csquotes}
\usepackage[left=2.5cm,right=2.2cm,top=2.5cm,bottom=3.5cm]{geometry}


\usepackage[
    backend=biber,
    style=alphabetic,
    sorting=ynt
]{biblatex}
\addbibresource{dep/term/reference.bib}

%%==================================================================
\newcommand\Z{\mathbb Z}
\let\le\leqslant
\let\ge\geqslant
\let\ph\varphi
\let\no\textnumero
%% и другие ваши макроопределения
%%==================================================================
\newtheorem{theorem}{Теорема}
\newtheorem{lemma}{Лемма}
\newtheorem{proposition}{Предложение}
\newtheorem{conj}{Гипотеза}
\theoremstyle{definition}
\newtheorem{definition}{Определение}
\DeclareMathOperator{\argmin}{argmin}
\DeclareMathOperator{\Det}{det}

%% и другие ваши типы теорем
%%==================================================================
\begin{document}
\thispagestyle{empty}
\begin{center}
ФЕДЕРАЛЬНОЕ ГОСУДАРСТВЕННОЕ БЮДЖЕТНОЕ\\
ОБРАЗОВАТЕЛЬНОЕ УЧРЕЖДЕНИЕ ВЫСШЕГО ОБРАЗОВАНИЯ\\
``САНКТ-ПЕТЕРБУРГСКИЙ ГОСУДАРСТВЕННЫЙ УНИВЕРСИТЕТ''\\
(СПбГУ)\\
\vfil
Образовательная программа бакалавриата ``Математика''
\vfil\vfil
\includegraphics[width=0.35\textwidth]{dep/spbu-logo.png}
\vfil\vfil
\textbf{\large Отчет о практике}\\
на тему\\
\textbf{Sequential модификация метода Flow Matching Posterior Estimation}
\end{center}

\vfil\vfil
\halign to \textwidth
{#\hfil\tabskip=0pt plus1fill &#\hfil\tabskip=0pt\cr
&
Выполнил студент 2 курса бакалавриата\cr
&
группа 24.Б03-мкн\cr
&
Ураков Михаил Алексеевич\cr
\noalign{\vfil}
&
Научный руководитель:\cr
&
Должность, ФИО\cr}

\vfil\vfil
\begin{center}
Санкт-Петербург\\
2026
\end{center}
\eject

\tableofcontents
\newpage

\section{Введение}
    Задача оценки истинных параметров математической модели встречается во многих научных и прикладных областях: 
    генетике \cite{abcinpopulationgenetics}, физике \cite{dax2023flowmatchingscalablesimulationbased}, финансах и др. 
    Она становится особенно сложной, когда невозможно аналитически описать распределение наблюдений 
    при заданных параметрах. В этой работе мы, следуя статье \cite{dax2023flowmatchingscalablesimulationbased}, 
    предлагаем модификацию метода Flow Matching Posterior Estimation и сравниваем его с другими методами на простых задачах, 
    а в конце применяем его для оценки параметеров модели стохастической волатильности.

    % Итак, наша обычно задача состоит из следующих частей
    % \begin{enumerate}
    %     \item Параметрическая модель некоторого процесса
    %     \item Симулятор 
    %     \item Данные наблюдений/симуляций
    % \end{enumerate} 
    % Цель обычно состоит в том, чтобы по наблюдениям подобрать наиболее согласованные с данными параметры модели

    % В настоящее время классичесие методы используют нейросети для аппроксимации искомого распределения параметров.
    % На выходе мы получаем функцию, которая по наблюдениям выдает параметры.
    % Но зачастую нам требуется знать параметры только для одного наблюдения, а не для всех возможных.

    % % Эта задача называется Sequential Neural Posterior Estimation

    % В этой работе мы сосредоточимся на Sequential модификации метода FMPE [DINGO-source], 
    % который является state-of-the-art на данный момент.

    С математической точки зрения наша задача формулируется следующим образом. 
    Будем обозначать $\theta \in \mathbb{R}^m$ --- вектор параметров модели, а $x \in \mathbb{R}^n$ --- наблюдаемые/симулированные данные
    или, чаще, их малоразмерная аппроксимация, выжимка. Из-за стохастической природы моделей удобно считать как $x$, 
    так и $\theta$ случайными величинами. 
    
    Далее, нам известно, возможно неявно, \textit{правдоподобие} $p(x \mid \theta)$ - распределение наблюдений $x$, которые мы можем получить при условии, 
    что модель инициализирована параметрами $\theta$. Также мы имеем вольность задать \textit{априорное} распределение $p(\theta)$ 
    - наше общее представление о том, в каких пределах могут изменяться параметры модели. 

    Имея все это, мы хотим узнать \textit{апостериорное} распределение $p(\theta \mid x)$ --- распределенине параметров, 
    которые могли бы дать нам такое же наблюдение $x$. Тогда по формуле байеса имеем 
    $$
    p(\theta \mid x) = \frac{p(x \mid \theta) p(\theta)}{p(x)}
    $$
    
    \paragraph{\textbf{Markov Chain Monte Carlo}} Таким образом, коль скоро нам известны $p(x \mid \theta), p(\theta)$, можно считать, 
    что $p(\theta \mid x) \propto p(x \mid \theta) p(\theta)$. Если бы мы имели аналитическое выражение для последнего, 
    то могли бы сэмплировать из апостериорного распределения методами Markov Chain Monte Carlo (MCMC, \cite{Chib2004}). 
    Их суть заключается в том, чтобы построить марковскую цепь, стационарное распределение которой будет совпадать 
    с заданным. Тогда, пройдя по ней достаточное число шагов, мы получим выборку из нужного распределения.
    
    Помимо того, что нам нужно знать явное выражение для ненормированной плотности, 
    у этого метода есть другие минусы, такие как: вычислительная затратность, сложность диагностики сходимости,
    чувствительность к настройкам и проблемы в больших размерностях.
    
    \paragraph{\textbf{Approximate Bayesian Computation}} В случае когда мы не можем аналитически выразить правдоподобие, 
    приходится действовать другими методами. В литературе их объединяют под общим названием ABC (Approximate Bayesian Computation \cite{abcinpopulationgenetics}) методов. 
    Классические алгоритмы ABC основаны на предоположении о том, что если наблюдения похожи, то и соответствующие параметры тоже.
    Более формально, они аппроксимируют истинное апостериорное $p(\theta \mid x_o)$ с помощью $p(\theta \mid \lvert x_o - x \rvert < \varepsilon)$
    
    Самый простой метод, Rejection ABC, работает так:
    \begin{enumerate}
        \item Дано: наблюдение $x_o$, хотим получить выборку из $p(\theta \mid x_o)$
        \item Генерируем $\theta^* \sim p(\theta)$
        \item Симулируем $x \sim p(x \mid \theta^*)$
        \item Если $\lvert x_o - x \rvert < \varepsilon$, то добавляем $\theta^*$ в выборку, иначе повторить 2.
    \end{enumerate} 
    Из очевидных минусов --- это проклятие размерности и неэффективность вычислений. 
    Для вычисления апостериора только в одной точке $x_o$ нам требуется провести тысячи симуляций, 
    большая часть из которых будет отметена, причем с увеличением размености это число растет экспоненциально.

    \paragraph{\textbf{Neural Posterior Estimation}} В последние годы, с развитием машинного обучения, 
    пришли и новые методы оценки апостериора. В отличие от классических алгоритмов, где апостериор вычислялся 
    только эмпирически, методы NPE (Neural Posterior Estimation, \cite{zammitmangion2024neuralmethodsamortizedinference}, 
    \cite{radev2020bayesflowlearningcomplexstochastic}) позволяют аппроксимировать функцию
    плотности с помощью нейросети. Часто в качестве универсального аппроксиматора используют \textit{нормализующие потоки} (Normalizing Flows, \cite{radev2020bayesflowlearningcomplexstochastic}),
    хорошо показавшие себя в задачах физики гравитационных волн \cite{dax2023flowmatchingscalablesimulationbased}. Нормализующие потоки отображают 
    сэмплы из одного (обычно, гауссовского) распределения в другое распределение серией последовательных несложных преобразований, 
    параметризованых нейросетью. 
    
    Минус такого подхода в том, что эти преобразования должны быть обратимы 
    и иметь определенный вид чтобы сеть можно было эффективно обучать. К тому же, такая архитектура плохо масштабируема 
    для нужд больших и сложных экспериментов. 
    
    Однако, самый большой и главный плюс в сравнении с классическими методами --- это то, что теперь, у нас на выходе
    не просто выборка из нужного распределения, а функция, которая каждому наблюдению $x$ сопоставляет распределение $p(\theta \mid x)$.
    Таким образом, один раз обучив такую модель, мы получим $p(\theta \mid x)$ для всех возможных $x$, а не только одного 
    --- в литературе это называется амортизированным инференсом (Amortized Inference, \cite{zammitmangion2024neuralmethodsamortizedinference}). 
    
    В контексте отображения одного распределения в другое эта задача родственна задаче генеративного ИИ. 
    А в последние годы, в этой области доминируют диффузионные модели \cite{yang2025diffusionmodelscomprehensivesurvey}. Принцип их работы 
    основан на том же преобразовании одного распределения в другое, только теперь сэмплы двигаются по траекториям, 
    которые являются решениями СДУ, заданного векторным полем. И нейросетью здесь параметризовано это самое поле.
    
    Детерминированным аналогом диффузионных моделей являются Flow Matching модели \cite{lipman2024flowmatchingguidecode}. В данном случае, 
    вместо стохастического дифференциального уравнения используется обыкновенное. На них мы и сосредоточимся в этой работе.
    

    
    % \paragraph{Normalizing Flows} Нормализующий поток (Normalizing Flow, \cite{NormFlow}) задает условное распределение параметров $q(\theta \mid x), \theta \in \mathbb{R}^n, \theta \in \mathbb{R}^m$, 
    % которое определяется как 
    % \[
    % q(\theta \mid x) = [\psi_x]_# q_0(\theta) = q_0(\psi_x^{-1}(\theta)) \det \Big| \partial_\theta \psi_x^{-1}(\theta) \Big|,
    % \]
    % где $\psi_x: \mathbf{R}^n \to \mathbb{R}^n$

    
    % \paragraph{Flow Matching}
    

\section{Flow Matching}
    Рассмотрим два распределения в $\mathbb R^d$: простое исходное
$p_0$, из которого легко генерировать сэмплы, и целевое $p_1$, которое
необходимо моделировать. Flow-модель преобразует случайную величину
$X_0\sim p_0$ в $X_1\sim p_1$ с помощью гладкого обратимого отображения
$\psi_1$. Если $\psi_1$ является диффеоморфизмом, плотность результата
определяется по формуле замены переменных:
\[
p_1(y) = (\psi_1)_\#p_0(y)
= p_0\bigl(\psi_1^{-1}(y)\bigr)
\left|\det \partial\psi_1^{-1}(y)\right|.
\]
Вместо того чтобы непосредственно строить такое преобразование,
рассмотрим его непрерывную во времени деформацию: для каждой начальной точки
$x$ траектория $\psi_t(x)$ удовлетворяет уравнению
\[
\frac{d}{dt}\psi_t(x) = v_t\bigl(\psi_t(x)\bigr),
\qquad \psi_0(x)=x,\qquad t\in[0,1].
\]
Если решение существует вплоть до $t=1$, то конечное отображение $\psi_1$ получается
интегрированием:
\[
\psi_1(x) = x + \int_0^{1} v_t\bigl(\psi_t(x)\bigr)\,dt.
\]
Таким образом, вместо поиска прямого преобразования достаточно задать
векторное поле $v_t$, которое определяет направление движения в каждой точке.
Это поле и становится объектом обучения в Flow Matching.

\begin{theorem}[Локальное существование и единственность потока, \cite{lipman2024flowmatchingguidecode}]
Пусть $v\in C^r([0,1]\times\mathbb R^d,\mathbb R^d)$ при $r\geq 1$.
Тогда задача Коши
\[
\dot\psi_t(x)=v_t(\psi_t(x)),\qquad \psi_0(x)=x
\]
имеет единственное локальное решение для каждого начального значения.
На областях, где эти решения определены, отображения между начальными и
текущими состояниями являются диффеоморфизмами.
\end{theorem}

Теорема гарантирует локальное решение, но сама по себе не утверждает, что
оно существует до $t=1$ для всех начальных точек. Для глобального решения
нужны дополнительные условия на поле; например, достаточным условием
является глобальная липшицевость по пространственной переменной на
рассматриваемом временном интервале. Поэтому далее утверждения о переносе
между $p_0$ и $p_1$ понимаются при условии, что соответствующий поток
определён на требуемом интервале времени.

\begin{proposition}
    Если $\psi_t$ удовлетворяет
    \[
    \frac{d}{dt}\psi_t(x) = v_t(\psi_t(x)),
    \]
    то плотность $p_t = [\psi_t]_\# p_0 (x)$ удовлетворяет уравнению непрерывности:
    \[
    \partial_t p_t(x) + \nabla \cdot \big(p_t(x) v_t(x)\big) = 0.
    \]
\end{proposition}
Тогда, если мы найдём правильное $v_t$, то $\psi_t$ \textit{хорошим образом} перенесёт $p_0$ в $p_1$
\begin{proof}
    Поскольку $\rho_t (y)$ определено как 
    \[
    \rho_t (y) = (\psi_t)_\# \rho_0 (y) = \rho_0(\psi_t^{-1}(y)) | \Det \partial_y \psi_t^{-1}(y)|
    \]
    то 
    \[
    \int \varphi(y)\,\rho_t(y)\,dy = \int \varphi(\psi_t(x))\,\rho_0(x)\,dx \quad \forall \varphi \in C_c^\infty.
    \]
    Продифференцируем обе части по $t$:
    \[
    \frac{d}{dt} \int \varphi(\psi_t(x))\,\rho_0(x)\,dx = \int \frac{d}{dt}\varphi(\psi_t(x))\,\rho_0(x)\,dx
    \]
    \[
    \frac{d}{dt}\varphi(\psi_t(x)) = \nabla \varphi(\psi_t(x)) \cdot \dot{\psi}_t(x).
    \]
    И т.к.
    \[
    \dot{\psi}_t(x) = u_t(\psi_t(x)),
    \]
    получим
    \[
    \frac{d}{dt} \int \varphi(\psi_t(x))\,\rho_0(x)\,dx = \int \nabla \varphi(\psi_t(x)) \cdot u_t(\psi_t(x))\,\rho_0(x)\,dx
    \]
    Снова, по определению $\rho_t (y)$
    \[
    \int \nabla \varphi(\psi_t(x)) \cdot u_t(\psi_t(x))\,\rho_0(x)\,dx = \int \nabla \varphi(y) \cdot u_t(y)\,\rho_t(y)\,dy
    \]
    Интегрированием по частям получим
    \[
    \int \nabla \varphi \cdot u_t \rho_t = - \int \varphi \, \nabla \cdot (\rho_t u_t)
    \]
    Тогда
    \[
    \frac{d}{dt} \int \varphi \rho_t + \int \varphi \, \nabla \cdot (\rho_t u_t) = 0
    \]
    И посольку это выполнено для всех $\varphi$, то
    \[
    \partial_t \rho_t + \nabla \cdot (\rho_t u_t) = 0
    \]
    в слабом смысле
\end{proof}

\subsection{Ключевая идея Flow Matching}
Основная идея Flow Matching --- задать путь распределений $(p_t)_{t\in[0,1]}$,
который соединяет исходное распределение $p_0$ с целевым $p_1$, а затем
обучить поле скоростей, порождающее этот путь. Если истинное поле $u_t$
известно, то его можно аппроксимировать нейронной сетью $v_t^\theta$ с помощью
обычной регрессии:
\[
\mathcal L_{\text{FM}}(\theta) =
\mathbb E_{t\sim U[0,1],\,X_t\sim p_t}
\left\|v_t^\theta(X_t)-u_t(X_t)\right\|^2.
\]
В этой постановке регрессия обучает не плотность $p_1$ непосредственно, а
локальное направление движения, необходимое для переноса вероятностной
массы вдоль пути. После обучения новые сэмплы получаются так: генерируем
$X_0\sim p_0$ и генерируем ОДУ с полем $v_t^\theta$ от $t=0$ до $t=1$.

Трудность состоит в том, что маргинальное поле $u_t(x)$ обычно нельзя
вычислить явно: оно зависит от всего распределения промежуточных состояний
$p_t$. Условная формулировка ниже заменяет эту недоступную цель на
аналитически известные условные поля, сохраняя тот же градиент обучения.

\subsection{Условная формулировка}
Пусть $Z\sim\pi(z)$ --- вспомогательная случайная величина, задающая
условный путь. Для каждого её значения $z$ определены условная плотность
$p_{t\mid z}(x\mid z)$ и поле $u_t(x\mid z)$, которое переносит эту
условную плотность во времени. Например, $z$ может задавать пару из начальной
и конечной точки $(x_0,x_1)$ и тем самым конкретную траекторию между ними.
Маргинальный путь получается усреднением условных плотностей:
\[
p_t(x)=\int p_{t\mid z}(x\mid z)\,\pi(dz).
\]
Соответствующее маргинальное поле является условным средним полем по
скрытой переменной при фиксированном текущем состоянии:
\[
u_t(x)=
\mathbb E\!\left[u_t(X_t\mid Z)\mid X_t=x\right],
\qquad X_t\sim p_t.
\]
И вместо него сеть обучают на условной регрессии:
\[
\mathcal L_{\text{cond}}(\theta)=
\mathbb E_{\substack{t\sim U[0,1],\,Z\sim\pi\\
X_t\sim p_{t\mid Z}(\cdot\mid Z)}}
\left\|v_t^\theta(X_t)-u_t(X_t\mid Z)\right\|^2.
\]
Для вычисления этого лосса достаточно сэмплировать время $t$ и условие
$z$, затем получить $X_t$ из условного пути и вычислить известную
условную скорость $u_t(X_t\mid z)$. Итого, мы пришли к Conditional Flow Matching, где обучение
--- это обычная задача регрессии по симулируемым парам "состояние — скорость".
Cвязь такого лосса с исходным маргинальным формализуется ниже.

\subsection{Пример}
\hfill \break
\hfill \break
Продемонстрируем на аналитическом примере упомянутые выше объекты. 
Разберем простейший случай с двумя гауссианами
\begin{align*}
    p_0(x) &= \mathcal N(\mu_0, \sigma_0) \\
    p_1(x) &= \mathcal N(\mu_1, \sigma_1) \\
    \psi_{t \mid 1}(x \mid x_1) &= t x_1 + (1 - t) x
\end{align*}
Тогда условное поле скоростей:
    \[
    u_{t \mid 1}(x \mid x_1) = \dot{\psi}_{t \mid 1} \circ \psi_{t \mid 1}^{-1}(x \mid x_1) = \frac{x_1 - x}{1 - t}
    \]
И условный probability path:
    \[
    p_{t \mid 1} (x \mid x_1) = p_{0 \mid 1} (\psi_{t \mid 1}^{-1}(x \mid x_1)) | \partial_x\psi_{t \mid 1}^{-1}(x \mid x_1) | = 
    \]
    \[
    = \frac{1}{\sqrt{2 \pi} \sigma} e^{-\frac{1}{2} \Big( \frac{\frac{x - t x_1}{1 - t} - \mu_0}{\sigma_0} \Big)^2} \cdot \frac{1}{1 - t} = 
    \mathcal N(t x_1 + (1 - t)\mu_0, (1 - t)\sigma_0)
    \]
Тогда marginal probability path:
    \[
    p_t(x) = \int p_{t \mid 1}(x | x_1) p_1(x_1) dx_1 = 
    \]
    \[
    = \int \mathcal N(t x_1 + (1 - t)\mu_0, (1 - t)\sigma_0) \> \mathcal N(\mu_1, \sigma_1) dx_1 =
    \]
    \[
    = \mathcal{N}(t \mu_1 + (1 - t) \mu_0, t^2 \sigma_1 + (1 - t) \sigma_0)
    \]
Отсюда по формуле Байеса 
    \[
    p_{1 \mid t}(x_1 \mid x) = \frac{p_{t \mid 1}(x \mid x_1) p_1(x_1)}{p_t(x)} = \mathcal N(\mu_{1 | t}(x), \sigma_{1|t}(x))
    \]
где
    \[
    \mu_{1 \mid t}(x) = \mu_1 + \frac{t \sigma_1 (x - t\mu_1 - (1 - t)\mu_0)}{t^2 \sigma_1 + (1 - t)\sigma_0}
    \]
    \[
    \sigma_{1 \mid t}(x) = \frac{\sigma_1 (1 - t) \sigma_0}{t^2 \sigma_1 + (1 - t)\sigma_0}
    \]
Итого, marginal velocity field:
    \[
    u_t(x) = \int u_{t \mid 1}(x \mid x_1) p_{1 \mid t}(x_1 \mid x) dx_1 = 
    \]
    \[
    = \mathbb{E}\Big[\frac{X_1 - X_t}{1 - t} \mid X_t = x\Big] = \frac{1}{1 - t}\mathbb{E}[X_1 \mid X_t = x] - \frac{x}{1 - t} =
    \]
    \[
    = \frac{\mu_{1 | t}(x) - x}{1 - t}
    \]

\subsection{Теорема о совпадении лоссов}
Условный лосс удобен вычислительно, но на первый взгляд отличается от
маргинального: в нем поле скоростей сети сопоставляется с $u_t(X_t\mid Z)$, а не с
$u_t(X_t)$. Ключевое свойство состоит в том, что эти две функции потерь имеют одинаковый градиент
по параметрам модели. При этом, конечно, значения самих лоссов могут отличаться на
слагаемое, не зависящее от параметров.

\begin{theorem}[\cite{lipman2024flowmatchingguidecode}]
Пусть $u_t(x)=\mathbb E[u_t(X_t\mid Z)\mid X_t=x]$, где совместное
распределение $(Z,X_t)$ задаётся как
$Z\sim\pi$, $X_t\mid Z\sim p_{t\mid Z}$. Тогда для дифференцируемого
поля $v_t^\theta$ и конечных ожиданий градиенты условного и маргинального
лоссов совпадают:
\[
\nabla_\theta \mathcal L_{\text{cond}}
=
\nabla_\theta \mathcal L_{\text{marg}},
\]
где
\[
\mathcal L_{\text{cond}}(\theta)=
\mathbb E_{t,Z,X_t}
\left\|v_t^\theta(X_t)-u_t(X_t\mid Z)\right\|^2,
\]
\[
\mathcal L_{\text{marg}}(\theta)=
\mathbb E_{t,X_t\sim p_t}
\left\|v_t^\theta(X_t)-u_t(X_t)\right\|^2.
\]
\end{theorem}

\begin{proof}
Обозначим $U=u_t(X_t\mid Z)$ и $m=u_t(X_t)$. По определению
маргинального поля $m=\mathbb E[U\mid t,X_t]$. Для фиксированных $t$ и
$X_t$ разложение квадратичного отклонения относительно условного среднего
даёт
\begin{align*}
&\mathbb E\!\left[
\left\|v_t^\theta(X_t)-U\right\|^2
\middle|\,t,X_t\right]\\
&\qquad =
\left\|v_t^\theta(X_t)-m\right\|^2 + 
\mathbb E\!\left[\left\|U-m\right\|^2
\middle|\,t,X_t\right].
\end{align*}
После усреднения по $t$ и $X_t$ получаем
$\mathcal L_{\mathrm{cond}}(\theta)
=\mathcal L_{\mathrm{marg}}(\theta)+C$, где
$C=\mathbb E\|U-m\|^2$ не зависит от $\theta$. Следовательно, градиенты двух лоссов совпадают.
\end{proof}


\section{Последовательные методы}
    В последовательных методах задача состоит в том, чтобы по одному
наблюдению $x_o$ эффективно оценить апостериорное распределение
$p(\theta\mid x_o)$. Вместо генерации всех обучающих симуляций из априорного
распределения алгоритм на каждом раунде использует proposal
$\tilde p_r(\theta)$, сосредоточенный в области параметров, которая
представляется перспективной. Из него генерируются пары
$\theta_i\sim\tilde p_r(\theta)$, $x_i\sim p(x\mid\theta_i)$, после чего
условная модель $q_\phi(\theta\mid x)$ обучается на накопленных данных.
Такой подход позволяет направлять вычислительный бюджет на параметры,
которые согласуются с наблюдением, но требует учитывать, что proposal
может отличаться от исходного априорного распределения.

В SMC-ABC апостериор приближается облаком частиц: параметры симулируются,
а частицы с симулированными данными, недостаточно близкими к наблюдению,
отбрасываются. Между популяциями порог расстояния уменьшается, а
оставшиеся частицы переиспользуются для предложения новых параметров.
Метод не обучает условную плотность и поэтому подходит как
симуляционный baseline, но требует выбрать расстояние и пороги принятия,
и один из очевидных минусов состоит в том, что число симуляций может
быстро расти, если вероятность принятия мала.

\subsection{Смещение при сэмплировании}

При обычном обучении максимизацией правдоподобия модель аппроксимирует
условное распределение, соответствующее proposal, из которого были
получены обучающие параметры:
\[
\tilde p(\theta\mid x)
=\frac{\tilde p(\theta)p(x\mid\theta)}{\tilde p(x)}.
\]
Искомое распределение определяется исходным априорным распределением:
\[
p(\theta\mid x)
=\frac{p(\theta)p(x\mid\theta)}{p(x)}
=\tilde p(\theta\mid x)
\frac{p(\theta)}{\tilde p(\theta)}
\frac{\tilde p(x)}{p(x)}.
\]
При фиксированном $x$ последний множитель не зависит от $\theta$ и
поглощается нормировкой, однако отношение $p(\theta)/\tilde p(\theta)$
в общем случае зависит от параметра. Поэтому без коррекции обучение на
данных из произвольного proposal приводит к смещению. Коррекция может
быть включена в функцию потерь, применена к обученной модели либо
устранена выбором proposal с подходящей структурой.

\subsection{Коррекция с помощью плотностей: SNPE-A и SNPE-B}

В SNPE-A \cite{papamakarios2016fast} смещение устраняется после обучения: сначала оценивается
$\tilde p(\theta\mid x)$, а затем результат умножается на отношение
$p(\theta)/\tilde p(\theta)$. В классическом варианте используются
смеси гауссовских распределений, для которых такое преобразование можно
выполнить аналитически. Это делает поправку явной, но связывает метод
с конкретным семейством моделей и требует подходящего вида как для
proposal, так и для апостериорной аппроксимации.

SNPE-B \cite{lueckmann2017flexible} переносит эту поправку в функцию потерь, используя веса
$w(\theta)=p(\theta)/\tilde p(\theta)$:
\[
\mathcal L_B(\phi)=
-\mathbb E_{\theta\sim\tilde p(\theta),\,x\sim p(x\mid\theta)}
\left[
\frac{p(\theta)}{\tilde p(\theta)}
\log q_\phi(\theta\mid x)
\right].
\]
В отличие от SNPE-A, этот подход не требует гауссовской смеси, но
нуждается в вычислении плотностей prior и proposal. Если proposal
назначает малую плотность параметрам, которые имеют заметную вероятность
при prior, веса могут иметь высокую дисперсию и затруднять обучение.

\subsection{SNPE-C: атомарная апостериорная оценка}

SNPE-C, также называемый APT \cite{greenberg2019apt}, на практике работает подобно классификатору.
Для каждого симулированного наблюдения $x_k$ строится набор
кандидатов параметров $\theta_1,\ldots,\theta_K$, и модель должна
назначить наибольшую оценку параметру, который породил это наблюдение.
В атомарном варианте оценки кандидатов пропорциональны
$q_\phi(\theta_j\mid x_k)/p(\theta_j)$:
\[
\mathcal L_C(\phi)=
-\mathbb E\left[
\log
\frac{q_\phi(\theta_k\mid x_k)/p(\theta_k)}
{\displaystyle\sum_{j=1}^{K}
q_\phi(\theta_j\mid x_k)/p(\theta_j)}
\right].
\]
Нормировка по конечному набору превращает оценку в задачу выбора
параметра, согласующегося с наблюдением. В библиотеке \texttt{sbi} этот подход
не требует явно вычислять плотность proposal и соответствующие веса,
поэтому удобен для последовательного обучения.
Однако лосс требует значений $\log q_\phi(\theta_j\mid x_k)$
для кандидатов и значения плотности prior. Таким образом, SNPE-C
устраняет необходимость в плотности proposal, но не необходимость в
оценке плотности апостериорной модели.

\subsection{Усечение prior}

Усечение proposal также используется в TSNPE \cite{deistler2022tsnpe}.
В Truncated SFMPE после первого раунда мы сужаем prior по области
$\mathcal S_r$, построенной из posterior предыдущего раунда. Для каждой
координаты параметра по сэмплам находим интервал, в который попадает
заданная доля значений, и расширяем его на величину \texttt{padding}.
В следующем раунде генерируем параметры из исходного prior и оставляем
те, что попали в эту область и допустимую область параметров. Поэтому
новый proposal — это prior,
ограниченный на $\mathcal S_r$:
\[
\tilde p(\theta)
=p(\theta\mid\theta\in\mathcal S_r)
=\frac{p(\theta)\mathbf 1_{\mathcal S_r}(\theta)}
{P(\theta\in\mathcal S_r)}.
\]

Мы не стали строить область по квантилям $\log q_\phi$: для FMPE вычисление
лог-плотности требует дополнительного решения ОДУ и оценки дивергенции
поля скоростей, что дорого при большом числе сэмплов. Вместо этого границы
задаются непосредственно по сэмплам posterior. Такой способ проще, но
область зависит от качества предыдущей оценки posterior; если prior редко
попадает в неё, отбраковка может потребовать много симуляций.


\section{Основные результаты}
    В этом разделе мы сравним SFMPE, SNPE-C и SMC-ABC на линейной гауссовской
модели, SIR и модели стохастической волатильности. Для каждого задания
рассматриваются бюджеты в 100, 1000 и 10000 симуляций на одно целевое
наблюдение; в каждой группе использовано 10 целевых наблюдений.
Для SFMPE и SNPE-C ниже приводятся результаты одного раунда, а SMC-ABC
является одношаговым методом. При сравнении качества posterior-выборок
используются 1000 сэмплов на целевое наблюдение.

\subsection{Линейная Gaussian-модель}

Пусть параметры $\theta\in\mathbb{R}^2$ имеют гауссовское априорное
распределение, а наблюдение $x\in\mathbb{R}^2$ задаётся линейной моделью
\[
\theta \sim \mathcal{N}(\mu_0,\Sigma_0), \qquad
x\mid\theta \sim \mathcal{N}(\theta+\delta,\Sigma_x).
\]
В эксперименте $\mu_0=(0,0)^\top$, $\delta=(0,0)^\top$,
\[
\Sigma_0 =
\begin{pmatrix}
1.5 & 0.3\\
0.3 & 1.0
\end{pmatrix},
\qquad
\Sigma_x =
\begin{pmatrix}
0.4 & 0.1\\
0.1 & 0.6
\end{pmatrix}.
\]
В качестве наблюдения используется сам вектор $x$, без дополнительного
сжатия. Для этой модели posterior известен точно:
\[
\theta\mid x \sim \mathcal{N}(\mu_{\mathrm{post}},\Sigma_{\mathrm{post}}),
\]
\[
\Sigma_{\mathrm{post}}
=\left(\Sigma_0^{-1}+\Sigma_x^{-1}\right)^{-1},
\qquad
\mu_{\mathrm{post}}
=\Sigma_{\mathrm{post}}\left(
\Sigma_0^{-1}\mu_0+\Sigma_x^{-1}(x-\delta)
\right).
\]
Поэтому здесь можно непосредственно сравнить выборку каждого алгоритма с
эталонным posterior с помощью C2ST: значение около $0.5$ соответствует
случайному угадыванию классификатора и, следовательно, неразличимости
двух выборок по этому тесту.

\subsection{SIR}
\hfill \break
\hfill \break
Здесь мы рассмотрим модель распространения заболевания SIR (Suspectable-Infected-Recovered)
Она устроена следующим образом. В каждый момент времени состояние описывается тремя числами $S, I, R$ 
- которые означают число людей из соответстующей категории ($S$ - Здоровые, $I$ - Больные, $R$ - Переболевшие). 
Всегда выполнено соотношение $S + I + R = N$ - общее число людей, а динамика задается так
\[
\begin{cases}
S_{t+1} = S_t - \Delta_{inf} \\
I_{t+1} = I_t + \Delta_{inf} - \Delta_{rec} \\
R_{t+1} = R_t + \Delta_{rec} \\
\end{cases}
\]
где $\Delta_{inf} \sim Binomial(S, 1 - e^{-\beta \frac{I}{N}}), \Delta_{rec} \sim Binomial(I, 1 - e^{-\gamma})$. 
Здесь $\beta, \gamma$ --- искомые параметры $\theta \in \mathbb{R}^2$, а $x \in \mathbb{R}^5$ --- вектор саммари-статистик, 
полученный из временного ряда $(S_t, I_t, R_t) \in \mathbb{R}^{T \times 3}$ 


\subsection{AR(1)-модель стохастической волатильности}

В модели стохастической волатильности наблюдаемые доходности $y_t$
зависят от скрытой логарифмической волатильности $v_t$:
\[
\begin{aligned}
v_0 &\sim \mathcal{N}\left(m,\frac{\sigma^2}{1-\phi^2}\right),\\
y_t &= \mu+\exp(v_{t-1}/2)\,\epsilon_t,\\
v_t &= m+\phi(v_{t-1}-m)+\sigma\eta_t,\qquad t=1,\ldots,T,
\end{aligned}
\]
где $\epsilon_t,\eta_t\sim\mathcal{N}(0,1)$ независимы. Параметр
$\mu$ задаёт средний уровень доходности, $\phi$ --- персистентность
лог-волатильности, $\sigma$ --- масштаб её инноваций, а $m$ --- её
долгосрочный средний уровень. Условие стационарности
$-1<\phi<1$ также является частью области допустимых параметров.

В бенчмарке используются независимые априорные распределения
\[
\mu\sim\mathcal{N}(0,0.2^2),\qquad
\phi=2U-1,\quad U\sim\operatorname{Beta}(20,1.5),
\]
\[
\sigma^2\sim\operatorname{InvGamma}(2.5,0.025),\qquad
m\sim\mathcal{N}(-1.5,1).
\]
Длина ряда $T$ случайно выбирается из целых чисел от 120 до 239. В
качестве наблюдений используются не исходные ряды, а пятикомпонентная
сводка
\[
x(y)=\left(
\bar y,
\log\widehat{\operatorname{Var}}(y),\
\widehat{\operatorname{skew}}(y),\
\widehat{\operatorname{kurt}}_{\mathrm{excess}}(y),\
\widehat{\rho}_1(y)
\right),
\]
где $\widehat{\rho}_1$ --- эмпирическая автокорреляция первого порядка.
Следовательно, оценивается распределение параметров при фиксированной
сводке ряда, а не при полном наблюдаемом временном ряде.

\subsection{Результаты бенчмарка}

В таблицах приведены среднее значение метрики и стандартное отклонение
по 10 целевым наблюдениям. В Gaussian-модели C2ST сравнивает posterior-
выборки с точным решением:
\[
\begin{array}{lccc}
\hline
\text{Метод} & 100 & 1000 & 10000\\
\hline
\text{SFMPE, 1 раунд} & 0.642\pm0.085 & 0.542\pm0.010 & 0.519\pm0.012\\
\text{SNPE-C, 1 раунд} & 0.561\pm0.021 & 0.524\pm0.007 & 0.501\pm0.010\\
\text{SMC-ABC} & 0.617\pm0.049 & 0.625\pm0.065 & 0.557\pm0.022\\
\hline
\end{array}
\]
SNPE-C ближе всего к эталонному posterior на всех трёх бюджетах; при
10000 симуляций значение $0.501$ практически совпадает с идеальным
$0.5$. У SFMPE также наблюдается улучшение при увеличении бюджета:
C2ST снижается с $0.642$ до $0.519$, однако при 10000 симуляций результат
остаётся несколько хуже SNPE-C. При бюджете 100 SMC-ABC (0.617) также
оказывается ближе к эталону, чем SFMPE (0.642); при бюджетах 1000 и
10000 SMC-ABC уступает обоим нейросетевым методам. Его качество при этом
не улучшается монотонно от 100 к 1000 симуляциям. При бюджете 100
разброс результатов SFMPE заметно больше, поэтому сравнение на этой
точке менее устойчиво.

Для StochVol качество оценивается по калибровке TARP. Используемая
здесь площадь отклонения от кривой ожидаемого покрытия (ATC) равна нулю
для идеальной калибровки; положительное значение соответствует
избыточному покрытию, а отрицательное --- недостаточному:
\[
\begin{array}{lccc}
\hline
\text{Метод} & 100 & 1000 & 10000\\
\hline
\text{SFMPE, 1 раунд} & 0.002 & 0.003 & 0.072\\
\text{SNPE-C, 1 раунд} & 0.074 & -0.058 & -0.028\\
\text{SMC-ABC} & 0.164 & 0.101 & 0.082\\
\hline
\end{array}
\]
На StochVol ATC у SFMPE почти равен нулю при бюджетах 100 и 1000,
однако увеличивается до $0.072$ при 10000 симуляциях. SNPE-C при
бюджете 100 имеет положительное отклонение от идеальной калибровки, а
при больших бюджетах --- отрицательное. У SMC-ABC ATC положителен на
всех трёх бюджетах и постепенно уменьшается с их ростом.

Для сравнения, на SIR-модели значения ATC для тех же методов составляют:
\[
\begin{array}{lccc}
\hline
\text{Метод} & 100 & 1000 & 10000\\
\hline
\text{SFMPE, 1 раунд} & 0.034 & 0.067 & -0.105\\
\text{SNPE-C, 1 раунд} & 0.129 & 0.091 & 0.006\\
\text{SMC-ABC} & 0.038 & 0.085 & 0.118\\
\hline
\end{array}
\]
На SIR значения ATC у SFMPE близки к нулю при бюджетах 100 и 1000, но
при 10000 становятся отрицательными. У SNPE-C ATC снижается с $0.129$
до $0.006$, а у SMC-ABC остаётся положительным и возрастает с бюджетом.

Для SFMPE и SNPE-C результаты при нескольких последовательных раундах
не дают одинакового эффекта на разных задачах и бюджетах. Например, на
StochVol при бюджете 1000 увеличение числа раундов SFMPE с одного до
трёх повышает ATC с $0.003$ до $0.141$.


\section{Заключение}
    В этой работе реализован вариант FMPE с последовательным уточнением
априорного распределения: на каждом раунде оно усекалось по
posterior-выборкам, полученным на предыдущем.

Эксперименты не показали, что дополнительные раунды систематически
повышают качество; в некоторых конфигурациях работа алгоритма становилась
нестабильной. На гауссовской задаче при одном раунде средний C2ST SNPE-C
был ближе к идеальному значению, чем у выбранной конфигурации SFMPE.

На SIR и StochVol соотношение методов по ATC менялось вместе с бюджетом
симуляций. Эти результаты следует трактовать с учётом малого числа
целевых наблюдений и того, что часть posterior-сэмплов оказывалась вне
допустимой области параметров. Поэтому имеющиеся эксперименты не дают
оснований говорить о безусловном превосходстве предложенного метода.


\section{Послесловие}
    Полученные выводы следует интерпретировать с учётом ограничений
бенчмарка. Для каждой комбинации задачи, метода и бюджета использовано
только 10 целевых наблюдений. Ограничение связано с вычислительной
стоимостью экспериментов; увеличение числа posterior-сэмплов для каждого
наблюдения не заменяет увеличение числа независимых целевых наблюдений.
Поэтому агрегированные показатели калибровки, включая TARP и SBC, имеют
ограниченную статистическую мощность и носят ориентировочный характер.

Для SIR и StochVol нет аналитического эталонного posterior, поэтому
качество оценивалось по калибровке и posterior-predictive проверкам.
Хорошая калибровка сама по себе не означает, что оценённая плотность
совпадает с истинным posterior. Кроме того, в StochVol в бенчмарке
используется пятикомпонентная сводка доходностей вместо полного временного
ряда; выводы относятся к posterior параметров при условии этой сводки.

Отдельное ограничение связано с поддержкой распределений. Для части
сэмплов SFMPE и SMC-ABC в StochVol нарушаются условия $-1<\phi<1$ и
$\sigma>0$; следовательно, соответствующие оценки калибровки нужно
рассматривать вместе с долей сэмплов, попавших в допустимую область.
Для SFMPE эта доля составляет приблизительно $0.33$, $0.36$ и $0.77$
при бюджетах 100, 1000 и 10000 соответственно; для SMC-ABC --- $0.61$,
$0.76$ и $0.78$. Поэтому близкий к нулю ATC SFMPE на малых бюджетах
нельзя интерпретировать отдельно от доли сэмплов вне допустимой области.

Наконец, измеренное время выполнения зависит от аппаратного устройства
и фоновой нагрузки, а запуски бенчмарка выполнялись с автоматическим
выбором устройства. Поэтому приведённые значения времени дают лишь
ориентировочное сравнение. Бюджет симуляций учитывает симуляции,
использованные для инференса, но не posterior-сэмплирование и
дополнительные симуляции для диагностик.

\section{Литература}
\printbibliography[heading=none]

% \begin{thebibliography}{9}%% Параметр определяет ширину номера, разумные значения при нумерации числами: 9 или 99
% \bibitem{OLD}
% C. L. Assic, \emph{The main properties of sipulki}, Main J.~Math. \textbf{1} (1900), \no 1, 1--100.

% \bibitem{NEW}
% С. О. Временник, \emph{Новые потрясающие свойства сипулек},
%   Российский мат. ж. \textbf{666} (2021), \no 1, 998--1000.

% \bibitem{EXTRA}
% И. И. Сидоров, \emph{Большая книга о сипульках}, изд-во Мир, Москва, 1962.
% \end{thebibliography}

\end{document}
%%==================================================================